\documentclass[a4paper,10pt,twoside]{article}
\usepackage{fontspec}
\setmainfont{Liberation Sans}
\setsansfont{Liberation Sans}
\renewcommand{\familydefault}{\sfdefault}
\usepackage{amsmath,amssymb}
\usepackage{geometry}
\usepackage{multicol}
\usepackage{graphicx}
\usepackage{fancyhdr}
\usepackage{lastpage}
\usepackage{booktabs}
\usepackage{listings}
\usepackage{eso-pic}
\usepackage{tikz}
\usepackage{array}
\usepackage{tabularx}
\usepackage{enumitem}
\usepackage{microtype}

\geometry{
    a4paper,
    top=1.8cm,
    bottom=1.8cm,
    left=1.2cm,
    right=1.2cm,
    headsep=14pt,
    footskip=20pt
}

\setlength{\columnsep}{16pt}
\setlength{\columnseprule}{0.35pt}

% Persistent Subtle Watermark on Every Page
\newsavebox{\watermarkbox}
\sbox{\watermarkbox}{%
    \begin{tabular}{c}
        \includegraphics[width=6.0cm]{/home/johan/GATE/GATE_CS_PYQ_DB/assets/studiva-butterfly-icon.jpg} \\[-0.2em]
        {\fontsize{14}{17}\selectfont\bfseries Studiva}
    \end{tabular}%
}
\AddToShipoutPictureBG{%
    \begin{tikzpicture}[remember picture, overlay]
        \node[opacity=0.035] at (current page.center) {\usebox{\watermarkbox}};
    \end{tikzpicture}%
}

% Running Headers & Footers
\pagestyle{fancy}
\fancyhf{}
\fancyhead[C]{\small\textbf{Studiva GATE Question Bank}}
\fancyhead[LE]{\small\textit{Computer Science \& IT (2011--2025)}}
\fancyhead[RO]{\small\textit{Studiva}}
\fancyfoot[LE,RO]{\scriptsize\textbf{Studiva} $\cdot$ \texttt{www.studiva.co.in}}
\fancyfoot[C]{\small\textbf{Page \thepage\ of \pageref{LastPage}}}
\fancyfoot[RE,LO]{\scriptsize\texttt{contact@studiva.co.in}}
\renewcommand{\headrulewidth}{0.4pt}
\renewcommand{\footrulewidth}{0.4pt}

% Listings Style (Clean, professional)
\lstset{
    basicstyle=\scriptsize\ttfamily,
    numbers=none,
    frame=single,
    rulecolor=\color{black!30},
    tabsize=2,
    breaklines=true,
    showstringspaces=false
}

\begin{document}

% ==================== PAGE 1: FORMAL TITLE & COVER PAGE ====================
\thispagestyle{empty}
\vspace*{1.5cm}
\begin{center}
    \includegraphics[width=2.5cm]{/home/johan/GATE/GATE_CS_PYQ_DB/assets/studiva-butterfly-icon.jpg}\\[1.5em]
    {\fontsize{24}{28}\selectfont\textbf{GATE COMPUTER SCIENCE \& IT}}\\[0.5em]
    {\fontsize{16}{20}\selectfont\textbf{PREVIOUS YEARS QUESTION BANK (2011 -- 2025)}}\\[0.6em]
    {\normalsize Detailed Examination Papers with Complete Official Answer Keys}\\[2.0em]
    \rule{0.65\textwidth}{0.8pt}\\[2.0em]
    {\large\textbf{Studiva GATE Examination Series}}\\[0.4em]
    {\small Comprehensive Authentic Question Bank for GATE Computer Science \& Information Technology}\\[2.8cm]

    \begin{tabular}{|ll|}
        \hline
        \textbf{Subject:} & Computer Science and Information Technology (CS/IT) \\
        \textbf{Period Covered:} & 2011 to 2025 \\
        \textbf{Examination Papers:} & 1 Complete Verified Paper Sets \\
        \textbf{Contents:} & Examination Questions \& Official Final Answer Keys \\
        \textbf{Published by:} & Studiva (\texttt{www.studiva.co.in}) \\
        \hline
    \end{tabular}

    \vfill
    {\small\textbf{Studiva $\cdot$ \texttt{www.studiva.co.in} $\cdot$ \texttt{contact@studiva.co.in}}}
\end{center}

\newpage

% ==================== PAGE 2: PUBLICATION IMPRINT & NOTES ====================
\thispagestyle{empty}
\vspace*{0.6cm}
\begin{center}
    {\large\textbf{PUBLICATION INFORMATION \& COPYRIGHT NOTICE}}\\[0.4em]
    \rule{0.6\textwidth}{0.5pt}
\end{center}

\vspace{1.2em}
\noindent\begin{tabular}{ll}
\textbf{Title:} & Studiva GATE Question Bank: Computer Science \& IT (2011--2025) \\
\textbf{Publisher:} & Studiva \\
\textbf{Website:} & \texttt{www.studiva.co.in} \\
\textbf{Official Inquiries:} & \texttt{contact@studiva.co.in} \\
\textbf{Edition:} & First Edition (2026) \\
\textbf{Syllabus Scope:} & Official GATE Syllabus (General Aptitude and CS / IT) \\
\end{tabular}

\vspace{1.8em}
\noindent\fbox{\parbox{\dimexpr\linewidth-2\fboxsep-2\fboxrule}{%
\textbf{COPYRIGHT NOTICE}\\[0.4em]
Copyright \copyright\ 2026 Studiva. All Rights Reserved.\\[0.4em]
No part of this publication may be reproduced, stored in a retrieval system, or transmitted in any form or by any means, electronic, mechanical, photocopying, recording, or otherwise, without prior written permission of Studiva, except for individual academic study.
}}

\vspace{1.5em}
\noindent\fbox{\parbox{\dimexpr\linewidth-2\fboxsep-2\fboxrule}{%
\textbf{PUBLISHER'S NOTE: ARCHIVAL QUESTION FORMAT}\\[0.4em]
This compendium has been meticulously curated and verified against official GATE question papers and final answer keys.
\par\medskip\noindent
\textbf{Notice on Image-Based Questions:}\\
Questions containing diagrams, circuit schematics, graphs, or image-based figures from the original examination papers have been omitted in this text edition due to non-readable format in historical source archives. All text-based conceptual, analytical, mathematical, algorithmic, and numerical questions are fully included with verified official keys.
}}

\vspace{1.5em}
\begin{center}
\textbf{EXAMINATION STRUCTURE \& PRACTICE GUIDELINES}
\end{center}
\begin{enumerate}[leftmargin=*]
    \item \textbf{Sectional Mark Division:} Each paper comprises General Aptitude (GA) and Computer Science (CS). Questions are clearly grouped into 1-mark and 2-mark blocks at the head of each section.
    \item \textbf{Question Types:} Includes Multiple Choice Questions (MCQ), Multiple Select Questions (MSQ), and Numerical Answer Type (NAT).
    \item \textbf{Official Answer Keys:} For self-evaluation without distraction, complete official answer keys are printed at the conclusion of each respective examination paper set.
\end{enumerate}

\newpage

% ==================== PAGE 3: TABLE OF CONTENTS ====================
\begin{center}
    {\Large\textbf{TABLE OF CONTENTS}}\\[0.4em]
    \rule{0.5\textwidth}{0.6pt}
\end{center}
\vspace{1em}

\noindent\begin{tabularx}{\linewidth}{|X|l|c|r|}
\hline
\textbf{Examination Paper} & \textbf{Session / Shift} & \textbf{Total Questions} & \textbf{Content Type} \\
\hline

GATE 2011 Computer Science & Shift CS & 54 & Questions + Answer Key \\ \hline
\end{tabularx}

\vspace{1.5em}
\noindent{\footnotesize\textit{Note: All questions in this edition have been typeset from verified digital records and cross-referenced with official answer keys released by the organising IIT/IISc institutes.}}

\newpage

% ==================== PAPER: 2011_CS ====================
\begin{center}
{\Large\textbf{GATE 2011 EXAMINATION}}\\[3pt]
{\normalsize\textbf{COMPUTER SCIENCE AND INFORMATION TECHNOLOGY}}\\[2pt]
{\small\textbf{Session / Shift: CS}}\\[5pt]
\rule{\linewidth}{0.8pt}\vspace{-4pt}
{\footnotesize\textbf{Duration: 180 Minutes} \hfill \textbf{Total Marks: 100} \hfill \textbf{Questions: 54}}\\
\vspace{-4pt}\rule{\linewidth}{0.4pt}
\end{center}
\vspace{4pt}
\noindent\fbox{\parbox{\dimexpr\linewidth-2\fboxsep-2\fboxrule}{%
\footnotesize\textbf{Examination Instructions:}\\[1pt]
1. Question Types: Multiple Choice Questions (MCQ), Multiple Select Questions (MSQ), and Numerical Answer Type (NAT).\\
2. Negative Marking Scheme: For 1-mark MCQ, $1/3$ mark is deducted; for 2-mark MCQ, $2/3$ mark is deducted. There is no negative marking for MSQ and NAT questions.\\
3. Mark Division: Questions are grouped by mark allocation within each section as specified below. Full official answer key is provided at the end of this paper.
}}
\vspace{10pt}
\begin{multicols}{2}

\vspace{8pt}
\noindent\textbf{\large\underline{COMPUTER SCIENCE AND INFORMATION TECHNOLOGY}}\\[2pt]
\noindent{\footnotesize\textit{[Section Block: 1-mark questions]}}
\vspace{4pt}

\noindent\textbf{Q.1} \hfill {\footnotesize\textsf{[MCQ]}}
\par\vspace{2pt}
{\bfseries In a compiler, keywords of a language are recognized during \par}
\par\vspace{2pt}\noindent
\textbf{(A)}~parsing of the program\par\vspace{2pt}\noindent
\textbf{(B)}~the code generation\par\vspace{2pt}\noindent
\textbf{(C)}~the lexical analysis of the program\par\vspace{2pt}\noindent
\textbf{(D)}~dataflow analysis
\par\vspace{6pt}\hrule height 0.25pt\vspace{6pt}
\noindent\textbf{Q.1} \hfill {\footnotesize\textsf{[MCQ]}}
\par\vspace{2pt}
{\bfseries Ta a (0, =<) where c is a constant
wry (FA, \par}
\par\vspace{3pt}\noindent\textbf{Answer:} \rule{3.2cm}{0.4pt}
\par\vspace{6pt}\hrule height 0.25pt\vspace{6pt}
\noindent\textbf{Q.2} \hfill {\footnotesize\textsf{[MCQ]}}
\par\vspace{2pt}
{\bfseries A layer-4 firewall (a device that can look at all protocol headers up to the transport layer) CANNOT \par}
\par\vspace{2pt}\noindent
\textbf{(A)}~block entire HTTP traffic during 9:00PM and 5:00AM\par\vspace{2pt}\noindent
\textbf{(B)}~block all ICMP traffic\par\vspace{2pt}\noindent
\textbf{(C)}~stop incoming traffic from a specific IP address but allow outgoing traffic to the same IP \newline address\par\vspace{2pt}\noindent
\textbf{(D)}~block TCP traffic from a specific user on a multi-user system during 9:00PM and 5:00AM.
\par\vspace{6pt}\hrule height 0.25pt\vspace{6pt}
\noindent\textbf{Q.2} \hfill {\footnotesize\textsf{[MCQ]}}
\par\vspace{2pt}
{\bfseries Ts, Ay (Fucaca «) where c; and c? are constants
The database can be configured to do ordered indexing on A, or hashing on A, - Which of the
following statements is TRUE? \par}
\par\vspace{2pt}\noindent
\textbf{(A)}~Ordered indexing will always outperform hashing for both queries\par\vspace{2pt}\noindent
\textbf{(B)}~Hashing will always outperform ordered indexing for both queries\par\vspace{2pt}\noindent
\textbf{(C)}~Hashing will outperform ordered indexing on Q1, but not on Q2\par\vspace{2pt}\noindent
\textbf{(D)}~Hashing will outperform ordered indexing on Q2, but not on QI \newline  \newline 2011 cs
\par\vspace{6pt}\hrule height 0.25pt\vspace{6pt}
\noindent\textbf{Q.3} \hfill {\footnotesize\textsf{[MCQ]}}
\par\vspace{2pt}
{\bfseries If two fair coins are flipped and at least one of the outcomes is known to be a head, what is the
probability that both outcomes are heads? \par}
\par\vspace{2pt}\noindent
\begin{tabular*}{\columnwidth}{@{\extracolsep{\fill}}p{0.46\columnwidth}p{0.46\columnwidth}@{}}
\textbf{(A)}~1/3 & \textbf{(B)}~1/4 \\
\addlinespace[2pt]
\textbf{(C)}~1/2 & \textbf{(D)}~2/3
\end{tabular*}
\par\vspace{6pt}\hrule height 0.25pt\vspace{6pt}
\noindent\textbf{Q.4} \hfill {\footnotesize\textsf{[MCQ]}}
\par\vspace{2pt}
{\bfseries Consider different activities related to email.
ml: Send an email from a mail client to a mail server
m2: Download an email from mailbox server to a mail client
m3: Checking email in a web browser
Which is the application level protocol used in each activity? \par}
\par\vspace{2pt}\noindent
\textbf{(A)}~ml: HTTP m2:SMTP\_ m3: POP\par\vspace{2pt}\noindent
\textbf{(B)}~ml: SMTP m2: FTP m3: HTTP\par\vspace{2pt}\noindent
\textbf{(C)}~ml: SMTP m2: POP m3: HTTP\par\vspace{2pt}\noindent
\textbf{(D)}~ml: POP m2: SMTP m3: IMAP
\par\vspace{6pt}\hrule height 0.25pt\vspace{6pt}
\noindent\textbf{Q.6} \hfill {\footnotesize\textsf{[MCQ]}}
\par\vspace{2pt}
{\bfseries Let the time taken to switch between user and kernel modes of execution be t; while the time taken
to switch between two processes be tf. Which of the following is TRUE? \par}
\par\vspace{2pt}\noindent
\textbf{(A)}~1) >t \newline B)t=h \newline (O)ti<t\par\vspace{2pt}\noindent
\textbf{(D)}~nothing can be said about the relation between tf, and f2 \newline CS-A 220 \newline 2011 cs
\par\vspace{6pt}\hrule height 0.25pt\vspace{6pt}
\noindent\textbf{Q.7} \hfill {\footnotesize\textsf{[MCQ]}}
\par\vspace{2pt}
{\bfseries A company needs to develop digital signal processing software for one of its newest inventions. The
software is expected to have 40000 lines of code. The company needs to determine the effort in
person-months needed to develop this software using the basic COCOMO model. The multiplicative
factor for this model is given as 2.8 for the software development on embedded systems, while the
exponentiation factor is given as 1.20. What is the estimated effort in person -months? \par}
\par\vspace{2pt}\noindent
\begin{tabular*}{\columnwidth}{@{\extracolsep{\fill}}p{0.46\columnwidth}p{0.46\columnwidth}@{}}
\textbf{(A)}~234.25 & \textbf{(B)}~932.50 \\
\addlinespace[2pt]
\textbf{(C)}~287.80 & \textbf{(D)}~122.40
\end{tabular*}
\par\vspace{6pt}\hrule height 0.25pt\vspace{6pt}
\noindent\textbf{Q.8} \hfill {\footnotesize\textsf{[MCQ]}}
\par\vspace{2pt}
{\bfseries Which of the following pairs have DIFFERENT expressive power? \par}
\par\vspace{2pt}\noindent
\textbf{(A)}~Embed web objects from different sites into the same page\par\vspace{2pt}\noindent
\textbf{(B)}~Refresh the page automatically after a specified interval\par\vspace{2pt}\noindent
\textbf{(C)}~Automatically redirect to another page upon download\par\vspace{2pt}\noindent
\textbf{(D)}~Display the client time as part of the page
\par\vspace{6pt}\hrule height 0.25pt\vspace{6pt}
\noindent\textbf{Q.10} \hfill {\footnotesize\textsf{[MCQ]}}
\par\vspace{2pt}
{\bfseries Which one of the following is NOT desired in a good Software Requirement Specifications (SRS)
document? \par}
\par\vspace{2pt}\noindent
\textbf{(A)}~Functional Requirements\par\vspace{2pt}\noindent
\textbf{(B)}~Non-Functional Requirements\par\vspace{2pt}\noindent
\textbf{(C)}~Goals of Implementation\par\vspace{2pt}\noindent
\textbf{(D)}~Algorithms for Software Implementation
\par\vspace{6pt}\hrule height 0.25pt\vspace{6pt}
\noindent\textbf{Q.11} \hfill {\footnotesize\textsf{[MCQ]}}
\par\vspace{2pt}
{\bfseries Acomputer handles several interrupt sources of which the following are relevant for this question.
¢ Interrupt from CPU temperature sensor (raises interrupt if CPU temperature is too high)
¢ Interrupt from Mouse (raises interrupt if the mouse is moved or a button is pressed)
© Interrupt from Keyboard (raises interrupt when a key is pressed or released)
¢ Interrupt from Hard Disk (raises interrupt when a disk read is completed)
Which one of these will be handled at the HIGHEST priority? \par}
\par\vspace{2pt}\noindent
\textbf{(A)}~Interrupt from Hard Disk\par\vspace{2pt}\noindent
\textbf{(B)}~Interrupt from Mouse\par\vspace{2pt}\noindent
\textbf{(C)}~Interrupt from Keyboard\par\vspace{2pt}\noindent
\textbf{(D)}~Interrupt from CPU temperature sensor \newline CS-A 3/20 \newline 2011 cs
\par\vspace{6pt}\hrule height 0.25pt\vspace{6pt}
\noindent\textbf{Q.12} \hfill {\footnotesize\textsf{[MCQ]}}
\par\vspace{2pt}
{\bfseries Consider a relational table with a single record for each registered student with the following
attributes.
1. Registration\_Num: Unique registration number of each registered student
2. UID: Unique identity number, unique at the national level for each citizen
3. BankAccount\_Num: Unique account number at the bank. A student can have multiple
accounts or joint accounts. This attribute stores the primary account number.
4. Name: Name of the student
5. Hostel\_Room: Room number of the hostel
Which of the following options is INCORRECT? \par}
\par\vspace{2pt}\noindent
\textbf{(A)}~(B) \newline (©)\par\vspace{2pt}\noindent
\textbf{(B)}~Registration\_Num can be a primary key\par\vspace{2pt}\noindent
\textbf{(C)}~UID is a candidate key if all students are from the same country\par\vspace{2pt}\noindent
\textbf{(D)}~p>
\par\vspace{6pt}\hrule height 0.25pt\vspace{6pt}
\noindent\textbf{Q.14} \hfill {\footnotesize\textsf{[MCQ]}}
\par\vspace{2pt}
{\bfseries The simplified SOP (Sum of Product) form of the Boolean expression
(P+O+R).(P+O+R).(P+O+R) is \par}
\par\vspace{2pt}\noindent
\textbf{(A)}~(P.O+R) ®) (P+OR) \newline © (P.O+R)\par\vspace{2pt}\noindent
\textbf{(D)}~(P.Q+R)
\par\vspace{6pt}\hrule height 0.25pt\vspace{6pt}
\noindent\textbf{Q.15} \hfill {\footnotesize\textsf{[MCQ]}}
\par\vspace{2pt}
{\bfseries The minimum number of D flip-flops needed to design a mod-258 counter is \par}
\par\vspace{2pt}\noindent
\begin{tabular*}{\columnwidth}{@{\extracolsep{\fill}}p{0.46\columnwidth}p{0.46\columnwidth}@{}}
\textbf{(A)}~9 & \textbf{(B)}~8 \\
\addlinespace[2pt]
\textbf{(C)}~512 & \textbf{(D)}~258
\end{tabular*}
\par\vspace{6pt}\hrule height 0.25pt\vspace{6pt}
\noindent\textbf{Q.16} \hfill {\footnotesize\textsf{[MCQ]}}
\par\vspace{2pt}
{\bfseries A thread is usually defined as a ``light weight process'' because an operating system (OS) maintains
smaller data structures for a thread than for a process. In relation to this, which of the following is
TRUE? \par}
\par\vspace{2pt}\noindent
\textbf{(A)}~On per-thread basis, the OS maintains only CPU register state\par\vspace{2pt}\noindent
\textbf{(B)}~The OS does not maintain a separate stack for each thread\par\vspace{2pt}\noindent
\textbf{(C)}~On per-thread basis, the OS does not maintain virtual memory state\par\vspace{2pt}\noindent
\textbf{(D)}~On per-thread basis, the OS maintains only scheduling and accounting information \newline CS-A 4/20 \newline 2011 cs
\par\vspace{6pt}\hrule height 0.25pt\vspace{6pt}
\noindent\textbf{Q.17} \hfill {\footnotesize\textsf{[MCQ]}}
\par\vspace{2pt}
{\bfseries K4 and Q3 are graphs with the following structures.
i
Which one of the following statements is TRUE in relation to these graphs? \par}
\par\vspace{2pt}\noindent
\textbf{(A)}~K4 is planar while Q3 is not\par\vspace{2pt}\noindent
\textbf{(B)}~Both K4 and Q3 are planar\par\vspace{2pt}\noindent
\textbf{(C)}~Q3 is planar while K3 is not\par\vspace{2pt}\noindent
\textbf{(D)}~Neither K4 nor Q3 is planar
\par\vspace{6pt}\hrule height 0.25pt\vspace{6pt}
\noindent\textbf{Q.18} \hfill {\footnotesize\textsf{[MCQ]}}
\par\vspace{2pt}
{\bfseries If the difference between the expectation of the square of a random variable (E [x >|) ana the
square of the expectation of the random variable (E [ xX ly is denoted by R , then \par}
\par\vspace{2pt}\noindent
\begin{tabular*}{\columnwidth}{@{\extracolsep{\fill}}p{0.46\columnwidth}p{0.46\columnwidth}@{}}
\textbf{(A)}~R=0 & \textbf{(B)}~R<0O \\
\addlinespace[2pt]
\textbf{(C)}~R20 & \textbf{(D)}~R>0
\end{tabular*}
\par\vspace{6pt}\hrule height 0.25pt\vspace{6pt}
\noindent\textbf{Q.19} \hfill {\footnotesize\textsf{[MCQ]}}
\par\vspace{2pt}
{\bfseries The lexical analysis for a modern computer language such as Java needs the power of which one of
the following machine models in a necessary and sufficient sense? \par}
\par\vspace{2pt}\noindent
\textbf{(A)}~Finite state automata\par\vspace{2pt}\noindent
\textbf{(B)}~Deterministic pushdown automata\par\vspace{2pt}\noindent
\textbf{(C)}~Non-deterministic pushdown automata\par\vspace{2pt}\noindent
\textbf{(D)}~Turing machine
\par\vspace{6pt}\hrule height 0.25pt\vspace{6pt}
\noindent\textbf{Q.20} \hfill {\footnotesize\textsf{[MCQ]}}
\par\vspace{2pt}
{\bfseries Let the page fault service time be 10 ms in a computer with average memory access time being
20 ns. If one page fault is generated for every 10° memory accesses, what is the effective access
time for the memory? \par}
\par\vspace{2pt}\noindent
\begin{tabular*}{\columnwidth}{@{\extracolsep{\fill}}p{0.46\columnwidth}p{0.46\columnwidth}@{}}
\textbf{(A)}~21 ns & \textbf{(B)}~30 ns \\
\addlinespace[2pt]
\textbf{(C)}~23 ns & \textbf{(D)}~35 ns
\end{tabular*}
\par\vspace{6pt}\hrule height 0.25pt\vspace{6pt}
\noindent\textbf{Q.21} \hfill {\footnotesize\textsf{[MCQ]}}
\par\vspace{2pt}
{\bfseries Consider a hypothetical processor with an instruction of type LW R1, 20(R2), which during
execution reads a 32-bit word from memory and stores it in a 32-bit register R1. The effective
address of the memory location is obtained by the addition of a constant 20 and the contents of
register R2. Which of the following best reflects the addressing mode implemented by this
instruction for the operand in memory? \par}
\par\vspace{2pt}\noindent
\textbf{(A)}~Immediate Addressing\par\vspace{2pt}\noindent
\textbf{(B)}~Register Addressing\par\vspace{2pt}\noindent
\textbf{(C)}~Register Indirect Scaled Addressing\par\vspace{2pt}\noindent
\textbf{(D)}~Base Indexed Addressing
\par\vspace{6pt}\hrule height 0.25pt\vspace{6pt}
\noindent\textbf{Q.22} \hfill {\footnotesize\textsf{[MCQ]}}
\par\vspace{2pt}
{\bfseries What does the following fragment of C program print?
Char c[] = ``'';
char *p = Cc;
printf (``\%s'', p + p[3] - pll]); \par}
\par\vspace{2pt}\noindent
\textbf{(A)}~(B) E2011\par\vspace{2pt}\noindent
\textbf{(C)}~2011\par\vspace{2pt}\noindent
\textbf{(D)}~011 \newline CS-A 5/20 \newline 2011 cs
\par\vspace{6pt}\hrule height 0.25pt\vspace{6pt}
\noindent\textbf{Q.24} \hfill {\footnotesize\textsf{[MCQ]}}
\par\vspace{2pt}
{\bfseries Let P bea regular language and Q be a context-free language such that Q Cc P.. (For example, let
P be the language represented by the regular expression p*g"and Q be \{p"q" |neé N\}). Then
which of the following is ALWAYS regular? \par}
\par\vspace{2pt}\noindent
\textbf{(A)}~PNO@\par\vspace{2pt}\noindent
\textbf{(B)}~P-Q\par\vspace{2pt}\noindent
\textbf{(C)}~X°-P () =*-Q
\par\vspace{6pt}\hrule height 0.25pt\vspace{6pt}
\noindent\textbf{Q.25} \hfill {\footnotesize\textsf{[MCQ]}}
\par\vspace{2pt}
{\bfseries An algorithm to find the length of the longest monotonically increasing sequence of numbers in an
array A[0:n-1] is given below.
Let L, denote the length of the longest monotonically increasing sequence starting at index iin
the array.
Initialize L,\_, =1.
For all i such thatO <i <n----2
L \{he if A[i] <A[i+1]
1 Otherwise
Finally the length of the longest monotonically increasing sequence is Max (Ly, Lye Ly ) :
Which of the following statements is TRUE? \par}
\par\vspace{2pt}\noindent
\textbf{(A)}~The aigorithm uses dynamic programming paradigm\par\vspace{2pt}\noindent
\textbf{(B)}~The algorithm has a linear complexity and uses branch and bound paradigm\par\vspace{2pt}\noindent
\textbf{(C)}~The algorithm has a non-linear polynomial complexity and uses branch and bound paradigm\par\vspace{2pt}\noindent
\textbf{(D)}~The algorithm uses divide and conquer paradigm. \newline CS-A 6/20 \newline 2011 cs \newline Q. 26 to Q. 55 carry two marks each.
\par\vspace{6pt}\hrule height 0.25pt\vspace{6pt}

\vspace{6pt}
\noindent{\footnotesize\textit{[Section Block: 2-mark questions]}}
\vspace{3pt}

\noindent\textbf{Q.26} \hfill {\footnotesize\textsf{[MCQ]}}
\par\vspace{2pt}
{\bfseries Consider the languages L1, L2 and L3 as given below.
L1=\{071"| p, gE N\},
L2=\{0'1"| p, ge Nand p=q\} and
L3=\{0°170" | p,q, r€N and p=q=r\}. Which of the following statements is NOT TRUE? \par}
\par\vspace{2pt}\noindent
\textbf{(A)}~Push Down Automata (PDA) can be used to recognize Li and L2\par\vspace{2pt}\noindent
\textbf{(B)}~L1 is a regular language\par\vspace{2pt}\noindent
\textbf{(C)}~All the three languages are context free\par\vspace{2pt}\noindent
\textbf{(D)}~Turing machines can be used to recognize all the languages
\par\vspace{6pt}\hrule height 0.25pt\vspace{6pt}
\noindent\textbf{Q.27} \hfill {\footnotesize\textsf{[MCQ]}}
\par\vspace{2pt}
{\bfseries Consider two binary operators `t' and `|' with the precedence of operator | being lower than that
of the operator t. Operator ¢ is right associative while operator | is left associative. Which one of
the following represents the parse tree for expression (713141312)? \par}
\par\vspace{2pt}\noindent
\textbf{(A)}~(B)\par\vspace{2pt}\noindent
\textbf{(C)}~(D) \newline CS-A 7120 \newline 2011 cs
\par\vspace{6pt}\hrule height 0.25pt\vspace{6pt}
\noindent\textbf{Q.28} \hfill {\footnotesize\textsf{[MCQ]}}
\par\vspace{2pt}
{\bfseries Ona non-pipelined sequential processor, a program segment, which is a part of the interrupt service
routine, is given to transfer 500 bytes from an I/O device to memory.
Initialize the address register
Initialize the count to 500
LOOP: Load a byte from device
Store in memory at address given by address register
Increment the address register
Decrement the count
If count $\ne$ 0 go to LOOP
Assume that each statement in this program is equivalent to a machine instruction which takes one
clock cycle to execute if it is a non-load/store instruction. The load-store instructions take two clock
cycles to execute.
The designer of the system also has an alternate approach of using the DMA controller to
implement the same transfer. The DMA controller requires 20 clock cycles for initialization and
other overheads. Each DMA transfer cycle takes two clock cycles to transfer one byte of data from
the device to the memory.
What is the approximate speedup when the DMA controller based design is used in place of the
interrupt driven program based input-output? \par}
\par\vspace{2pt}\noindent
\begin{tabular*}{\columnwidth}{@{\extracolsep{\fill}}p{0.46\columnwidth}p{0.46\columnwidth}@{}}
\textbf{(A)}~3.4 & \textbf{(B)}~4.4 \\
\addlinespace[2pt]
\textbf{(C)}~5.1 & \textbf{(D)}~6.7
\end{tabular*}
\par\vspace{6pt}\hrule height 0.25pt\vspace{6pt}
\noindent\textbf{Q.29} \hfill {\footnotesize\textsf{[MCQ]}}
\par\vspace{2pt}
{\bfseries We are given a set of distinct elements and an unlabeled binary tree with n nodes. In how many
ways can we populate the tree with the given set so that it becomes a binary search tree? \par}
\par\vspace{2pt}\noindent
\textbf{(A)}~0\par\vspace{2pt}\noindent
\textbf{(B)}~1\par\vspace{2pt}\noindent
\textbf{(C)}~nt\par\vspace{2pt}\noindent
\textbf{(D)}~y-"C, \newline nT
\par\vspace{6pt}\hrule height 0.25pt\vspace{6pt}
\noindent\textbf{Q.30} \hfill {\footnotesize\textsf{[MCQ]}}
\par\vspace{2pt}
{\bfseries Which one of the following options is CORRECT given three positive integers x, y and z, anda
predicate
P(x) A(x = 1) A Vy(Se(x = y*2>(y=x)v(\_y =1)) \par}
\par\vspace{2pt}\noindent
\textbf{(A)}~P(x) being true means that xis a prime number\par\vspace{2pt}\noindent
\textbf{(B)}~P(x)being true means that xis a number other than 1\par\vspace{2pt}\noindent
\textbf{(C)}~P(x)is always true irrespective of the value of x\par\vspace{2pt}\noindent
\textbf{(D)}~P(x) being true means that x has exactly two factors other than | and x \newline m2 woe \newline - .. . os cos x +1sin x
\par\vspace{6pt}\hrule height 0.25pt\vspace{6pt}
\noindent\textbf{Q.31} \hfill {\footnotesize\textsf{[MCQ]}}
\par\vspace{2pt}
{\bfseries Given i=-/\~{}1, what will be the evaluation of the definite integral J ------\_ dx?
9 COSx---isinx \par}
\par\vspace{2pt}\noindent
\textbf{(A)}~0\par\vspace{2pt}\noindent
\textbf{(B)}~2\par\vspace{2pt}\noindent
\textbf{(C)}~-i\par\vspace{2pt}\noindent
\textbf{(D)}~i \newline CS-A 8/20 \newline 2011 cs
\par\vspace{6pt}\hrule height 0.25pt\vspace{6pt}
\noindent\textbf{Q.32} \hfill {\footnotesize\textsf{[MCQ]}}
\par\vspace{2pt}
{\bfseries Consider a database table T containing two columns X and Y each of type integer. After the
creation of the table, one record (X=1, Y=1) is inserted in the table.
Let MX and MY denote the respective maximum values of X and Y among all records in the table
at any point in time. Using MX and MY, new records are inserted in the table 128 times with X and
Y values being MX+1, 2*MY+1 respectively. It may be noted that each time after the insertion,
values of MX and MY change.
What will be the output of the following SQL query after the steps mentioned above are carried
out?
SELECT Y FROM T WHERE X=7; \par}
\par\vspace{2pt}\noindent
\textbf{(A)}~127\par\vspace{2pt}\noindent
\textbf{(B)}~255 (©) 129\par\vspace{2pt}\noindent
\textbf{(D)}~257
\par\vspace{6pt}\hrule height 0.25pt\vspace{6pt}
\noindent\textbf{Q.33} \hfill {\footnotesize\textsf{[MCQ]}}
\par\vspace{2pt}
{\bfseries Consider a finite sequence of random values X = [X,,\%,,--.,\%, ]. Let M, be the mean and 0, be
the standard deviation of X. Let another finite sequence Y of equal length be derived from this
as y, =a*x, +b, where a and bare positive constants. Let 44, be the mean and oy, be the standard
deviation of this sequence. Which one of the following statements is INCORRECT? \par}
\par\vspace{2pt}\noindent
\textbf{(A)}~1/5\par\vspace{2pt}\noindent
\textbf{(B)}~4/25 (©) 1/4\par\vspace{2pt}\noindent
\textbf{(D)}~2/5
\par\vspace{6pt}\hrule height 0.25pt\vspace{6pt}
\noindent\textbf{Q.35} \hfill {\footnotesize\textsf{[MCQ]}}
\par\vspace{2pt}
{\bfseries Consider the following table of arrival time and burst time for three processes PO, P1 and P2.
Process | Arrival time | Burst Time
PO Oms 9 ms
Pl lms 4 ms
P2 2 ms 9 ms
The pre-emptive shortest job first scheduling algorithm is used. Scheduling is carried out only at
arrival or completion of processes. What is the average waiting time for the three processes? \par}
\par\vspace{2pt}\noindent
\textbf{(A)}~5.0 ms\par\vspace{2pt}\noindent
\textbf{(B)}~4.33 ms\par\vspace{2pt}\noindent
\textbf{(C)}~6.33 ms\par\vspace{2pt}\noindent
\textbf{(D)}~7.33 ms \newline CS-A 9/20 \newline 2011 \newline cs
\par\vspace{6pt}\hrule height 0.25pt\vspace{6pt}
\noindent\textbf{Q.39} \hfill {\footnotesize\textsf{[MCQ]}}
\par\vspace{2pt}
{\bfseries Consider evaluating the following expression tree on a machine with load-store architecture in
which memory can be accessed only through load and store instructions. The variables a, b, c,d
and e are initially stored in memory. The binary operators used in this expression tree can be
evaluated by the machine only when the operands are in registers. The instructions produce result
only in a register. If no intermediate results can be stored in memory,
what is the minimum number
of registers needed to evaluate this expression? \par}
\par\vspace{2pt}\noindent
\textbf{(A)}~248000\par\vspace{2pt}\noindent
\textbf{(B)}~44000\par\vspace{2pt}\noindent
\textbf{(C)}~19000\par\vspace{2pt}\noindent
\textbf{(D)}~25000 \newline Consider a relational table r with sufficient number of records, having attributes A, , A,,..A, and \newline let 1 <p <n. Two queries Q] and Q2 are given below.
\par\vspace{6pt}\hrule height 0.25pt\vspace{6pt}
\noindent\textbf{Q.40} \hfill {\footnotesize\textsf{[MCQ]}}
\par\vspace{2pt}
{\bfseries Consider the matrix as given below.
12 3
047
00 3
Which one of the following options provides the CORRECT values of the eigenvalues of the
matrix? \par}
\par\vspace{2pt}\noindent
\begin{tabular*}{\columnwidth}{@{\extracolsep{\fill}}p{0.46\columnwidth}p{0.46\columnwidth}@{}}
\textbf{(A)}~1, 4,3 & \textbf{(B)}~3, 7,3 \\
\addlinespace[2pt]
\textbf{(C)}~7, 3,2 & \textbf{(D)}~1, 2,3
\end{tabular*}
\par\vspace{6pt}\hrule height 0.25pt\vspace{6pt}
\noindent\textbf{Q.42} \hfill {\footnotesize\textsf{[MCQ]}}
\par\vspace{2pt}
{\bfseries Definition of a language L with alphabet \{a\} is given as following.
L =\{a™ |k >0, and n is a positive integer constant\}
What is the minimum number of states needed in a DFA to recognize L? \par}
\par\vspace{2pt}\noindent
\textbf{(A)}~k+1\par\vspace{2pt}\noindent
\textbf{(B)}~n+] («© 2"\par\vspace{2pt}\noindent
\textbf{(D)}~2"
\par\vspace{6pt}\hrule height 0.25pt\vspace{6pt}
\noindent\textbf{Q.43} \hfill {\footnotesize\textsf{[MCQ]}}
\par\vspace{2pt}
{\bfseries An 8KB direct-mapped write-back cache is organized as multiple blocks, each of size 32-bytes. The
processor generates 32-bit addresses. The cache controller maintains the tag information for each
cache block comprising of the following.
1 Valid bit
1 Modified bit
As many bits as the minimum needed to identify the memory block mapped in the cache.
What is the total size of memory needed at the cache controller to store meta-data (tags) for the
cache? \par}
\par\vspace{2pt}\noindent
\begin{tabular*}{\columnwidth}{@{\extracolsep{\fill}}p{0.46\columnwidth}p{0.46\columnwidth}@{}}
\textbf{(A)}~4864 bits & \textbf{(B)}~6144 bits \\
\addlinespace[2pt]
\textbf{(C)}~6656 bits & \textbf{(D)}~5376 bits
\end{tabular*}
\par\vspace{6pt}\hrule height 0.25pt\vspace{6pt}
\noindent\textbf{Q.44} \hfill {\footnotesize\textsf{[MCQ]}}
\par\vspace{2pt}
{\bfseries An application loads 100 libraries at startup. Loading each library requires exactly one disk access.
The seek time of the disk to a random location is given as 10 ms. Rotational speed of disk is
6000 rpm. If all 100 libraries are loaded from random locations on the disk, how long does it take to
load all libraries? (The time to transfer data from the disk block once the head has been positioned
at the start of the block may be neglected.) \par}
\par\vspace{2pt}\noindent
\textbf{(A)}~0.50 s\par\vspace{2pt}\noindent
\textbf{(B)}~1.50s\par\vspace{2pt}\noindent
\textbf{(C)}~1.258\par\vspace{2pt}\noindent
\textbf{(D)}~1.00 s \newline Cs-A 11/20 \newline 2011 \newline cs
\par\vspace{6pt}\hrule height 0.25pt\vspace{6pt}
\noindent\textbf{Q.46} \hfill {\footnotesize\textsf{[MCQ]}}
\par\vspace{2pt}
{\bfseries A deterministic finite automaton (DFA) D with alphabet ¥ = \{a,b\} is given below.
Which of the following finite state machines is a valid minimal DFA which accepts the same
language as D? \par}
\par\vspace{2pt}\noindent
\textbf{(A)}~3\par\vspace{2pt}\noindent
\textbf{(B)}~9\par\vspace{2pt}\noindent
\textbf{(C)}~5\par\vspace{2pt}\noindent
\textbf{(D)}~6 \newline  \newline 2011 cs
\par\vspace{6pt}\hrule height 0.25pt\vspace{6pt}
\noindent\textbf{Q.47} \hfill {\footnotesize\textsf{[MCQ]}}
\par\vspace{2pt}
{\bfseries The following is the comment written for a C function.
/* This function computes the roots of a quadratic equation
a.x*2 + b.x + c = 0. The function stores two real roots
in *rootl and *root2 and returns the status of validity
of roots. It handles four different kinds of cases.
(i) When coefficient a is zero irrespective of discriminant
(ii) When discriminant is positive
(iii) When discriminant is zero
(iv) When discriminant is negative.
Only in case (ii) and (iii), the stored roots are valid.
Otherwise 0 is stored in the roots. The function returns
0 when the roots are valid and -1 otherwise.
The function also ensures root1l $\ge$ root2.
int get\_QuadRoots(float a, float b, float c,
float *root1l, float *root2);
*/
A software test engineer is assigned the job of doing black box testing. He comes up with the
following test cases, many of which are redundant.
Test | Input Set Expected Output Set
Case | a b c\_\_| rootl | root2 | Return Value
TI | 0.0 0.0 |\_7.0 0.0 0.0 -1
T2 | 0.0 10| 3.0 0.0 0.0 -1
T3\_ | 1.0 2.0| 10] -10\} -1.0 0
T4 | 4.0 |-12.0| 9.0 1.5 1.5 1)
TS | 1.0 | -2.0 | -3.0 3.0 | -1.0 0
To | 1.0 10] 4.0 0.0 0.0 -1
Which one of the following options provide the set of non-redundant tests using equivalence class
partitioning approach from input perspective for black box testing? \par}
\par\vspace{2pt}\noindent
\textbf{(A)}~T1, T2, T3, T6\par\vspace{2pt}\noindent
\textbf{(B)}~T1, T3, T4, TS\par\vspace{2pt}\noindent
\textbf{(C)}~T2, T4, T5, T6\par\vspace{2pt}\noindent
\textbf{(D)}~T2, T3, T4, TS \newline CS-A 13/20 \newline 2011 \newline cs \newline Common Data Questions \newline Common Data for Questions 48 and 49: \newline Consider the following recursive C function that takes two arguments. \newline unsigned int foo(unsigned int n, unsigned int r) \{ \newline if (m>0) return ((n\%r) + foo(n/r, r)); \newline else return 0;
\par\vspace{6pt}\hrule height 0.25pt\vspace{6pt}
\noindent\textbf{Q.48} \hfill {\footnotesize\textsf{[MCQ]}}
\par\vspace{2pt}
{\bfseries --- What is the return value of the function £00 when it is called as £00 (345, 10)? \par}
\par\vspace{2pt}\noindent
\begin{tabular*}{\columnwidth}{@{\extracolsep{\fill}}p{0.46\columnwidth}p{0.46\columnwidth}@{}}
\textbf{(A)}~345 & \textbf{(B)}~12 \\
\addlinespace[2pt]
\textbf{(C)}~5 & \textbf{(D)}~3
\end{tabular*}
\par\vspace{6pt}\hrule height 0.25pt\vspace{6pt}
\noindent\textbf{Q.49} \hfill {\footnotesize\textsf{[MCQ]}}
\par\vspace{2pt}
{\bfseries --- What is the return value of the function £00 when it is called as foo(513, 2)? \par}
\par\vspace{2pt}\noindent
\textbf{(A)}~9\par\vspace{2pt}\noindent
\textbf{(B)}~8\par\vspace{2pt}\noindent
\textbf{(C)}~5\par\vspace{2pt}\noindent
\textbf{(D)}~2 \newline Common Data for Questions 50 and 51: \newline Consider the following circuit involving three D-type flip-flops used in a certain type of counter \newline configuration.
\par\vspace{6pt}\hrule height 0.25pt\vspace{6pt}
\noindent\textbf{Q.50} \hfill {\footnotesize\textsf{[MCQ]}}
\par\vspace{2pt}
{\bfseries If at some instance prior to the occurrence of the clock edge, P, Q and R have a value 0, | and 0
respectively, what shall be the value of PQR after the clock edge? \par}
\par\vspace{2pt}\noindent
\begin{tabular*}{\columnwidth}{@{\extracolsep{\fill}}p{0.46\columnwidth}p{0.46\columnwidth}@{}}
\textbf{(A)}~000 & \textbf{(B)}~001 \\
\addlinespace[2pt]
\textbf{(C)}~010 & \textbf{(D)}~011
\end{tabular*}
\par\vspace{6pt}\hrule height 0.25pt\vspace{6pt}
\noindent\textbf{Q.52} \hfill {\footnotesize\textsf{[MCQ]}}
\par\vspace{2pt}
{\bfseries --- Thecost of link N2-N3 reduces to 2 (in both directions). After the next round of updates, what will
be the new distance vector at node, N3? \par}
\par\vspace{2pt}\noindent
\begin{tabular*}{\columnwidth}{@{\extracolsep{\fill}}p{0.46\columnwidth}p{0.46\columnwidth}@{}}
\textbf{(A)}~(3, 2, 0, 2, 5) & \textbf{(B)}~(3, 2, 0, 2, 6) \\
\addlinespace[2pt]
\textbf{(C)}~(7, 2, 0, 2, 5) & \textbf{(D)}~(7, 2, 0, 2, 6)
\end{tabular*}
\par\vspace{6pt}\hrule height 0.25pt\vspace{6pt}
\noindent\textbf{Q.54} \hfill {\footnotesize\textsf{[MCQ]}}
\par\vspace{2pt}
{\bfseries What will be the cost of the minimum spanning tree (MST) of such a graph with n nodes? \par}
\par\vspace{2pt}\noindent
\begin{tabular*}{\columnwidth}{@{\extracolsep{\fill}}p{0.46\columnwidth}p{0.46\columnwidth}@{}}
\textbf{(A)}~bU1n?-5n) & \textbf{(B)}~rn? -nt1 \\
\addlinespace[2pt]
\textbf{(C)}~6n-11 & \textbf{(D)}~2n+1
\end{tabular*}
\par\vspace{6pt}\hrule height 0.25pt\vspace{6pt}
\noindent\textbf{Q.55} \hfill {\footnotesize\textsf{[MCQ]}}
\par\vspace{2pt}
{\bfseries The length of the path from v, to v, in the MST of previous question with n = 10 is \par}
\par\vspace{2pt}\noindent
\textbf{(A)}~11\par\vspace{2pt}\noindent
\textbf{(B)}~25\par\vspace{2pt}\noindent
\textbf{(C)}~31\par\vspace{2pt}\noindent
\textbf{(D)}~41 \newline CS-A 15/20 \newline 2011 cs \newline General Aptitude (GA) Questions
\par\vspace{6pt}\hrule height 0.25pt\vspace{6pt}
\noindent\textbf{Q.56} \hfill {\footnotesize\textsf{[MCQ]}}
\par\vspace{2pt}
{\bfseries Which of the following options is the closest in the meaning to the word below:
Inexplicable \par}
\par\vspace{2pt}\noindent
\begin{tabular*}{\columnwidth}{@{\extracolsep{\fill}}p{0.46\columnwidth}p{0.46\columnwidth}@{}}
\textbf{(A)}~Incomprehensible & \textbf{(B)}~Indelible \\
\addlinespace[2pt]
\textbf{(C)}~Inextricable & \textbf{(D)}~Infallible
\end{tabular*}
\par\vspace{6pt}\hrule height 0.25pt\vspace{6pt}
\noindent\textbf{Q.57} \hfill {\footnotesize\textsf{[MCQ]}}
\par\vspace{2pt}
{\bfseries If Log (P) = (1/2)Log (Q) = (1/3) Log (R), then which of the following options is TRUE? \par}
\par\vspace{2pt}\noindent
\textbf{(A)}~P?= QR?\par\vspace{2pt}\noindent
\textbf{(B)}~V=PR (©) V=R?P\par\vspace{2pt}\noindent
\textbf{(D)}~R=P°Q
\par\vspace{6pt}\hrule height 0.25pt\vspace{6pt}
\noindent\textbf{Q.58} \hfill {\footnotesize\textsf{[MCQ]}}
\par\vspace{2pt}
{\bfseries Choose the most appropriate word(s) from the options given below to complete the following
sentence.
I contemplated Singapore for my vacation but decided against it. \par}
\par\vspace{2pt}\noindent
\begin{tabular*}{\columnwidth}{@{\extracolsep{\fill}}p{0.46\columnwidth}p{0.46\columnwidth}@{}}
\textbf{(A)}~to visit & \textbf{(B)}~having to visit \\
\addlinespace[2pt]
\textbf{(C)}~visiting & \textbf{(D)}~for a visit
\end{tabular*}
\par\vspace{6pt}\hrule height 0.25pt\vspace{6pt}
\noindent\textbf{Q.59} \hfill {\footnotesize\textsf{[MCQ]}}
\par\vspace{2pt}
{\bfseries Choose the most appropriate word from the options given below to complete the following
sentence.
If you are trying to make a strong impression on your audience, you cannot do so by being
understated, tentative or . \par}
\par\vspace{2pt}\noindent
\begin{tabular*}{\columnwidth}{@{\extracolsep{\fill}}p{0.46\columnwidth}p{0.46\columnwidth}@{}}
\textbf{(A)}~hyperbolic & \textbf{(B)}~restrained \\
\addlinespace[2pt]
\textbf{(C)}~argumentative & \textbf{(D)}~indifferent
\end{tabular*}
\par\vspace{6pt}\hrule height 0.25pt\vspace{6pt}
\noindent\textbf{Q.60} \hfill {\footnotesize\textsf{[MCQ]}}
\par\vspace{2pt}
{\bfseries Choose the word from the options given below that is most nearly opposite in meaning to the given
word:
Amalgamate \par}
\par\vspace{2pt}\noindent
\textbf{(A)}~merge\par\vspace{2pt}\noindent
\textbf{(B)}~split\par\vspace{2pt}\noindent
\textbf{(C)}~collect\par\vspace{2pt}\noindent
\textbf{(D)}~separate \newline Q. 61 to Q. 65 carry two marks each.
\par\vspace{6pt}\hrule height 0.25pt\vspace{6pt}
\noindent\textbf{Q.61} \hfill {\footnotesize\textsf{[MCQ]}}
\par\vspace{2pt}
{\bfseries Few school curricula include a unit on how to deal with bereavement and grief, and yet all
students at some point in their lives suffer from losses through death and parting.
Based on the above passage which topic would not be included in a unit on bereavement? \par}
\par\vspace{2pt}\noindent
\textbf{(A)}~how to write a letter of condolence\par\vspace{2pt}\noindent
\textbf{(B)}~what emotional stages are passed through in the healing process\par\vspace{2pt}\noindent
\textbf{(C)}~what the leading causes of death are\par\vspace{2pt}\noindent
\textbf{(D)}~how to give support to a grieving friend \newline CS-A 16/20 \newline 2011 cs
\par\vspace{6pt}\hrule height 0.25pt\vspace{6pt}
\noindent\textbf{Q.63} \hfill {\footnotesize\textsf{[MCQ]}}
\par\vspace{2pt}
{\bfseries The variable cost (V) of manufacturing a product varies according to the equation V= 4q, where q
is the quantity produced. The fixed cost (F) of production of same product reduces with q according
to the equation F = 100/q. How many units should be produced to minimize the total cost (V+F)? \par}
\par\vspace{2pt}\noindent
\begin{tabular*}{\columnwidth}{@{\extracolsep{\fill}}p{0.46\columnwidth}p{0.46\columnwidth}@{}}
\textbf{(A)}~5 & \textbf{(B)}~4 \\
\addlinespace[2pt]
\textbf{(C)}~7 & \textbf{(D)}~6
\end{tabular*}
\par\vspace{6pt}\hrule height 0.25pt\vspace{6pt}
\noindent\textbf{Q.64} \hfill {\footnotesize\textsf{[MCQ]}}
\par\vspace{2pt}
{\bfseries A transporter receives the same number of orders each day. Currently, he has some pending orders
(backlog) to be shipped. If he uses 7 trucks, then at the end of the 4th day he can clear all the
orders. Alternatively, if he uses only 3 trucks, then all the orders are cleared at the end of the 10th
day. What is the minimum number of trucks required so that there will be no pending order at the
end of the Sth day? \par}
\par\vspace{2pt}\noindent
\begin{tabular*}{\columnwidth}{@{\extracolsep{\fill}}p{0.46\columnwidth}p{0.46\columnwidth}@{}}
\textbf{(A)}~4 & \textbf{(B)}~5 \\
\addlinespace[2pt]
\textbf{(C)}~6 & \textbf{(D)}~7
\end{tabular*}
\par\vspace{6pt}\hrule height 0.25pt\vspace{6pt}
\noindent\textbf{Q.65} \hfill {\footnotesize\textsf{[MCQ]}}
\par\vspace{2pt}
{\bfseries A container originally contains 10 litres of pure spirit. From this container 1 litre of spirit is
replaced with | litre of water. Subsequently, | litre of the mixture is again replaced with | litre of
water and this process is repeated one more time. How much spirit is now left in the container? \par}
\par\vspace{2pt}\noindent
\textbf{(A)}~7.58 litres\par\vspace{2pt}\noindent
\textbf{(B)}~7.84 litres\par\vspace{2pt}\noindent
\textbf{(C)}~7 litres\par\vspace{2pt}\noindent
\textbf{(D)}~7.29 litres \newline  \newline CS-A 17/20
\par\vspace{6pt}\hrule height 0.25pt\vspace{6pt}
\end{multicols}
\vspace{8pt}
\begin{center}
\rule{0.35\linewidth}{0.5pt}\\[2pt]
\textbf{\small --- END OF QUESTION PAPER ---}\\[-8pt]
\rule{0.35\linewidth}{0.5pt}
\end{center}
\vspace{10pt}
\subsection*{OFFICIAL ANSWER KEY --- GATE 2011 (CS)}
\vspace{-4pt}
\noindent{\footnotesize\textit{Official final answer keys published by the GATE Organising Institute. For NAT questions, allowable numerical ranges are displayed where applicable. MTA = Marks To All.}}
\vspace{4pt}
\begin{center}
\small
\resizebox{\linewidth}{!}{%
\begin{tabular}{|c|c|l||c|c|l||c|c|l||c|c|l|}
\hline
\textbf{Q} & \textbf{Type} & \textbf{Key} & \textbf{Q} & \textbf{Type} & \textbf{Key} & \textbf{Q} & \textbf{Type} & \textbf{Key} & \textbf{Q} & \textbf{Type} & \textbf{Key} \\
\hline\hline
1 & MCQ & C & 16 & MCQ & C & 31 & MCQ & D & 52 & MCQ & A \\ \hline
1 & MCQ & C & 17 & MCQ & B & 32 & MCQ & A & 54 & MCQ & B \\ \hline
2 & MCQ & D & 18 & MCQ & C & 33 & MCQ & D & 55 & MCQ & C \\ \hline
2 & MCQ & D & 19 & MCQ & A & 35 & MCQ & A & 56 & MCQ & A \\ \hline
3 & MCQ & A & 20 & MCQ & B & 39 & MCQ & C & 57 & MCQ & B \\ \hline
4 & MCQ & C & 21 & MCQ & D & 40 & MCQ & A & 58 & MCQ & C \\ \hline
6 & MCQ & C & 22 & MCQ & C & 42 & MCQ & B & 59 & MCQ & B \\ \hline
7 & MCQ & A & 24 & MCQ & C & 43 & MCQ & D & 60 & MCQ & D \\ \hline
8 & MCQ & B & 25 & MCQ & A & 44 & MCQ & B & 61 & MCQ & C \\ \hline
10 & MCQ & D & 26 & MCQ & C & 46 & MCQ & C & 63 & MCQ & A \\ \hline
11 & MCQ & D & 27 & MCQ & B & 47 & MCQ & C & 64 & MCQ & C \\ \hline
12 & MCQ & A & 28 & MCQ & A & 48 & MCQ & B & 65 & MCQ & D \\ \hline
14 & MCQ & B & 29 & MCQ & B & 49 & MCQ & D &  & &  \\ \hline
15 & MCQ & A & 30 & MCQ & A & 50 & MCQ & D &  & &  \\ \hline
\end{tabular}%
}
\end{center}
\end{document}